% !TEX root = ../main.tex
\chapter{便携式 EMT 系统硬件平台设计与实现}
\markboth{便携式 EMT 系统硬件平台设计与实现}{}

基于第二章对高导电率铁磁性目标 EMT 正逆问题的理论分析，以及第三、四章对 SCF-DeepONet 成像模型的设计与验证，本文进一步面向工程实现，设计并研制了一套便携式 EMT 硬件平台。与传统依赖台式信号源、锁相放大器、工控机和外置采集卡拼接而成的实验系统相比，道岔现场检测设备不仅需要具备稳定的激励能力和微弱信号提取能力，还必须兼顾体积、功耗、供电方式、抗干扰能力以及现场使用的便利性。

针对上述需求，本文采用``两板协同''的总体硬件架构：一块为信号切换板，主要负责阵列线圈的激励选通、接收通道复用以及激励电流监测；另一块为测控板，主要负责电源转换、激励信号产生与功率放大、微弱信号调理与锁相解调、模数转换以及主控通信等功能。系统以 $100\,\mathrm{kHz}$ 作为当前主工作点，配套 $4\times4$ 平面线圈阵列，实现对铁磁性试块表面及近表面伤损的电磁响应采集。同时，激励链路和参考链路预留了一定的频率调节能力，其实际可调范围主要受板级 PLL 锁定范围及模拟链路工作带宽限制，从而为后续多频 EMT 实验提供硬件基础。

本章将依次介绍便携式 EMT 系统的总体方案、平面传感器阵列与多路选通网络、激励信号产生与功率放大电路、微弱信号调理与双相锁相放大器、模数转换与无线通信模块、电源管理及启机时序，以及整机 PCB 抗干扰设计。

\section{系统总体方案设计与需求分析}

\subsection{便携式道岔检测性能需求分析}

用于道岔伤损检测的 EMT 系统，本质上是在高导电率、高磁导率金属背景中，探测局部缺陷引起的微弱电磁扰动。对于真实道岔钢轨或等效铁磁性试块而言，强背景场、强趋肤效应和复杂现场环境共同决定了该系统在硬件实现上具有较高要求。结合本课题的检测对象和工作模式，硬件平台需要满足以下几方面需求。

\textbf{(1) 感性负载驱动与激励稳定性需求。}

EMT 阵列线圈在 $100\,\mathrm{kHz}$ 工作频率下呈现明显感性特征。根据前期建模与实测结果，单个阵列线圈的阻抗模值约为 $195\,\Omega$，其虚部占主导。为在线圈中建立幅值稳定、相位可控的交变磁场，激励链路不仅需要提供足够的输出电压摆幅，还必须具备驱动感性负载的能力，避免在高频工作下出现明显波形畸变或相位漂移。

\textbf{(2) 微弱信号高信噪比提取需求。}

在强一次磁场背景下，伤损引起的二次感应电压通常远小于原场耦合信号，属于典型的窄带微弱交流信号测量问题。若直接采用高速 ADC 对原始 $100\,\mathrm{kHz}$ 波形进行宽带采样，不仅会增加系统功耗和数据吞吐压力，也会对 ADC 动态范围和前端噪声控制提出较高要求。因此，系统需要在模拟前端完成相敏检波，将交流响应转换为同相与正交两路低频或直流分量，再进行高精度数字化采集。这样既可提高信噪比和动态储备，也便于后续直接构造 EMT 成像和 SCF-DeepONet 推理所需的复数观测数据。

\textbf{(3) 阵列分时复用与通道隔离需求。}

本文系统采用 $4\times4$ 平面线圈阵列，并按照``单线圈激励、其余线圈接收''的方式工作。一帧完整测量需要依次完成 16 轮激励，并在每轮激励下采集其余 15 个线圈的响应。因此，硬件平台必须具备可靠的多路选通能力：一方面，激励信号需要被定向送入当前激励线圈；另一方面，其余线圈的接收信号需要被依次复用至统一采集链路。更重要的是，激励通道中的较大功率交流信号与接收通道中的微弱感应信号必须充分隔离，以降低大信号串扰对测量结果的影响。

\textbf{(4) 便携供电与高集成度需求。}

铁路道岔现场检测不适合长期依赖体积庞大的台式仪器组合，因此系统需要支持外部直流电源或电池供电，并在有限板级面积内集成电源管理、信号产生、功率放大、锁相解调、模数转换、主控与无线通信等功能模块。与此同时，为兼顾模拟测量精度和数字控制灵活性，系统还需要在高度集成的前提下合理划分功能区域，降低大信号、高频数字信号与微弱模拟信号之间的相互干扰。

\textbf{(5) 工作频率可调与参数扩展需求。}

本文当前硬件设计与实验验证主要围绕 $100\,\mathrm{kHz}$ 单频工作点展开。这一频率既便于与有限元仿真和 SCF-DeepONet 训练数据保持一致，也有利于验证整机信号链路的稳定性。与此同时，从 EMT 检测机理看，不同频率下的磁场穿透深度和伤损敏感性存在差异。因此，激励源与参考链路仍需具备一定频率调节能力，为后续多频实验和参数优化预留空间。需要说明的是，该能力主要受板级 PLL 的锁定范围、锁相放大器参考链路以及模拟滤波器带宽共同限制。

\textbf{(6) 现场防护与电磁兼容需求。}

道岔检测环境通常伴随振动、粉尘、静电积累以及复杂电磁干扰。由于系统前端包含高增益微弱信号调理电路，后端又包含 MCU、DDS、ADC 等低压精密器件，若外部静电放电或瞬态脉冲通过传感器接口、电源入口或通信接口耦合进入主板，可能引起测量漂移、逻辑异常、系统复位甚至器件损坏。因此，系统在硬件设计中需要充分考虑接口防护、电源滤波、接地策略、模拟数字隔离以及关键模拟链路的屏蔽设计，以提高现场运行可靠性和测量重复性。

\subsection{硬件系统总体架构设计}

基于上述需求分析，并综合考虑系统体积、功耗、抗干扰能力和调试便利性，本文最终采用``测控板 + 信号切换板''的双板式总体硬件架构。该结构一方面避免将大功率激励切换网络与高增益微弱信号调理电路完全混布在同一块板上，降低了串扰风险；另一方面也便于对激励链路、接收链路和选通网络进行独立调试与后续优化。

如图 \ref{fig:EMT_system} 所示，系统完整工作过程可概括为``激励产生—通道分发—电磁感应—微弱信号调理—锁相解调—高精度采集—无线传输''链路。

\begin{figure}[!htbp]
    \centering
    \includegraphics[width=0.82\textwidth]{figures/sys_structure.jpg}
    \figcaption{便携式 EMT 系统总体硬件架构框图}{Overall hardware architecture of the portable EMT system}
    \label{fig:EMT_system}
\end{figure}

在激励端，测控板首先由直接数字频率合成器产生 $100\,\mathrm{kHz}$ 正弦信号。该信号经过差分恢复、前置电压放大和功率放大后，形成具备负载驱动能力的交变激励源，并通过信号切换板定向送至当前选中的激励线圈。激励线圈在铁磁性试块附近建立一次交变磁场。当试块中存在伤损时，局部电导率与磁导率分布变化会改变涡流路径和磁场分布，其结果表现为其余接收线圈感应电压的幅值与相位变化。

在接收端，来自信号切换板的微弱感应电压首先进入高输入阻抗缓冲与程控增益调理链路，以实现前端隔离、阻抗匹配和动态范围扩展；随后进入双相锁相放大链路，在同源正交参考信号作用下被解调为同相与正交两路直流量。经电平位移与输入保护后，两路直流量由高分辨率 ADC 采集，并送入主控单元完成数据组织、校准管理和无线传输。

在控制与成像层面，主控单元负责通道切换、DDS 配置、增益控制、ADC 采样和数据封包；上位机则负责数据解析、差分归一化、传统 EMT 成像或 SCF-DeepONet 模型推理。由此，整套系统形成了从电磁激励、阵列采集到伤损检测结果输出的完整闭环。

该架构的核心优势在于：通过模拟前端的窄带相敏检波，显著降低了后端数字采样率和数据吞吐压力；通过双板式结构实现激励网络与微弱采集链路的功能分离，提升了便携式系统在复杂现场中的抗干扰能力和测量稳定性。

\section{传感器阵列与多路选通模块设计}

传感器阵列及其多路选通网络位于整个 EMT 硬件平台的最前端，直接决定了目标区域电磁扰动信息的获取质量。对于本文所设计的 $4\times4$ 阵列式 EMT 系统而言，前端设计不仅需要保证线圈参数一致性，还需要在大电流激励与微弱接收信号共存的条件下，实现可靠的分时复用与通道隔离。

\subsection{平面传感器阵列设计}

结合道岔相关构件的几何特征和有限元仿真结果，本文设计了一种 $4\times4$ 拓扑结构的 16 线圈平面传感器阵列。与传统环形阵列相比，平面阵列更适合贴近钢轨、试块或道岔表面布置，能够在有限检测面积内形成规则空间采样布局，并与后续有限元模型和深度学习数据组织方式保持一致。

阵列中每个探测线圈均由线径为 $0.3\,\mathrm{mm}$ 的铜制漆包线绕制，线圈匝数为 200 匝，骨架内径为 $4.5\,\mathrm{mm}$，外径为 $7\,\mathrm{mm}$，高度为 $8.4\,\mathrm{mm}$。接口方面，阵列采用 I-PEX 连接器实现与后端硬件的可插拔连接，以兼顾小型化、机械可靠性和现场维护便利性。为保证安装精度，本文配套设计了 3D 打印线圈底座和专用接口 PCB，确保线圈极性、排列顺序和磁场叠加方向与理论模型保持一致。

\begin{figure}[!htbp]
    \centering
    \includegraphics[width=0.62\textwidth]{figures/EMT_COILS.jpg}
    \figcaption{$4\times4$ 平面传感器阵列实物图}{Physical image of the $4\times4$ planar sensor array}
    \label{fig:EMT_COILS}
\end{figure}

对于基于相敏检波的单频 EMT 系统而言，线圈参数一致性直接影响激励稳定性和解调相位一致性。因此，本文利用 HP 4285A 精密阻抗分析仪，对 16 个线圈在不同频率下的复阻抗进行了逐一扫频测量。经过绕制工艺控制和手工圈数微调后，所有线圈在 $100\,\mathrm{kHz}$ 附近的阻抗模值均保持在约 $195\,\Omega$，最大值与最小值之差为 $0.268\,\Omega$，标准差为 $0.055$，表明阵列线圈具有较高的一致性。

\begin{figure}[!htb]
    \centering
    \includegraphics[width=0.72\textwidth]{figures/coil_impedance.jpg}
    \figcaption{传感器线圈阻抗特性测试实验照}{Impedance characterization of sensor coils}
    \label{fig:coil_impedance}
\end{figure}

此外，在保持线圈几何参数和工作频率一致的条件下，本文还利用有限元仿真对线圈阻抗与电磁场分布进行了数值验证。仿真结果与实测结果整体趋势一致，实测均值与仿真结果之间的相对误差约为 $0.211\%$。该结果说明，所设计的线圈阵列不仅在制造一致性方面满足要求，其电磁行为也与仿真模型具有较好一致性，为后续系统标定、硬件测试和算法验证提供了可靠基础。

\begin{figure}[!htb]
    \centering
    \includegraphics[width=0.76\textwidth]{figures/coil_consistency.jpg}
    \figcaption{传感器线圈阻抗一致性分析}{Impedance consistency analysis of sensor coils}
    \label{fig:coil_consistency}
\end{figure}

\subsection{收发异构大电流多路选通电路设计}

在 16 线圈 EMT 系统中，一帧完整测量需要依次将 16 个线圈轮流设为激励线圈，并在每一轮激励下采集其余线圈的接收响应。因此，多路选通网络不仅决定了阵列分时复用的实现方式，也直接影响激励链路与接收链路之间的隔离程度。

本文系统中，激励端工作在 $100\,\mathrm{kHz}$ 条件下，具有较大的电压摆幅和感性负载驱动特征；而接收端感应信号幅值较小，对附加相位误差、寄生电容和串扰极为敏感。若采用统一的对称式模拟开关架构，同时承担激励与接收通道切换，难以兼顾大电流耐受能力和微弱信号高保真传输需求。基于这一考虑，本文采用``机械继电器 + 低寄生参数模拟复用器''的收发异构多路选通架构：激励端使用机械继电器完成大电流选通，接收端使用低寄生参数模拟复用器完成微弱信号复用输出。

\textbf{1. 大电流激励信号选通设计}

针对激励端的大信号工况，本文选用 16 颗 G5V-1 DC5 单刀双掷微型机械继电器构成激励选通网络。继电器公共端连接对应阵列线圈，常开端连接激励总线，常闭端连接接收支路。当某一线圈被选为激励线圈时，该线圈自动脱离接收链路并接入激励通路，从硬件连接关系上降低了大电流激励对接收前端造成冲击的风险。

在控制逻辑上，16 路继电器由 4-16 译码器 CD74H154 统一译码选通，并通过后级驱动电路完成继电器控制。与纯模拟开关方案相比，机械继电器切换速度较慢，但在耐压、耐流和隔离能力方面更适合感性负载下的激励信号切换任务。

\begin{figure}[!htbp]
    \centering
    \includegraphics[width=0.80\textwidth]{figures/relay.jpg}
    \figcaption{信号切换控制电路原理图}{Schematic of relay-based signal switching control circuit}
    \label{fig:G5V-1}
\end{figure}

\textbf{2. 低寄生参数接收信号选通设计}

在接收支路，阵列线圈输出的感应信号经继电器常闭端接入多路复用网络。考虑到本系统采用 $100\,\mathrm{kHz}$ 单频相敏测量，接收信号对通道附加寄生参数较为敏感，因此本文选用 ADI 公司的 16 选 1 模拟多路复用器 ADG1206 作为接收端信号选通核心器件。

系统工作时，处于接收状态的线圈信号经各自继电器常闭端汇入 ADG1206 对应输入端，再由 ADG1206 公共端顺序输出至测控板的统一采集链路。该结构使激励端和接收端能够根据各自电气特性采用不同器件：继电器负责激励端硬隔离，ADG1206 负责接收端小信号复用。为进一步抑制高阻抗节点上的串扰与辐射耦合，PCB 布局中对接收支路采用短走线、包地和敏感节点隔离等措施。

\textbf{3. 非对称防护与激励电流监测设计}

由于激励端与接收端在信号幅度、阻抗水平和防护需求上存在显著差异，本文在信号切换板上采用非对称保护与监测设计。

在激励端，激励总线入口配置瞬态电压抑制器件，用于吸收外部尖峰脉冲。同时，为获得激励电流参考信息，系统在激励输入端串入 PA1005.100NLT 高频电流互感器及采样电阻，将激励电流转换为可测电压信号，并通过独立继电器在``接收信号测量''和``激励电流测量''之间切换。该设计为后续分析不同激励线圈下的实际注入一致性及测量归一化提供了硬件基础。

\begin{figure}[!htb]
    \centering
    \includegraphics[width=0.78\textwidth]{figures/PA1005.jpg}
    \figcaption{PA1005 电流互感器原理图}{Schematic of PA1005 current transformer}
    \label{fig:PA1005}
\end{figure}

在接收端，由于后级采集链路具有高输入阻抗，保护电路应尽量避免对正常微弱信号传输造成加载。基于这一原则，本文在接收支路采用``限流电阻 + 低漏电流钳位二极管''的保护方案，在异常过压时限制进入模拟复用器和后级放大器的电流，同时尽量减小正常测量条件下的信号衰减。

综上，本文所设计的收发异构多路选通架构，将激励链路的大信号耐受能力和接收链路的微弱信号传输需求进行解耦。机械继电器负责激励端硬隔离与定向选通，ADG1206 负责接收端小信号复用，电流互感器提供激励参考观测通道。该模块不仅满足了 16 线圈 EMT 系统分时复用测量需求，也为后续锁相解调和高精度采集链路提供了稳定前端接口。

\section{激励产生与功率放大模块设计}

激励链路负责在阵列线圈中建立稳定、可控的交变磁场，其性能直接影响一次场强度、接收响应幅值以及系统整体信噪比。对于本文设计的便携式 EMT 系统，激励模块需要在 $100\,\mathrm{kHz}$ 工作频率下输出稳定正弦信号，并具备驱动感性线圈负载的能力。基于这一需求，本文在测控板上采用``直接数字频率合成 + 差分恢复 + 前置电压放大 + 功率放大''的激励链路结构。

\subsection{激励信号发生电路}

本系统采用 ADI 公司的低功耗直接数字频率合成器 AD9834 作为激励信号源核心器件。AD9834 可通过 SPI 接口由主控 MCU 配置，其输出频率 $f_{out}$ 由主时钟频率 $f_{MCLK}$ 与 28 位频率控制字共同决定，满足
\begin{equation}
\label{eqn:DDS_freq}
f_{out}=\frac{\mathrm{FREQREG}\times f_{MCLK}}{2^{28}}.
\end{equation}

在本文系统中，AD9834 的主时钟频率设置为 $75\,\mathrm{MHz}$，当前主要围绕 $100\,\mathrm{kHz}$ 单频工作点展开设计与验证。利用 MCU 对频率控制字的配置，可在后续实验中实现一定范围内的频率调整。

由于 AD9834 输出为两路反相信号，后级电路需要先完成差分恢复和单端转换。本文采用精密差分放大器 AD8276 作为前级接口。AD8276 内部集成高匹配精度电阻网络，能够较好地完成差分转单端处理，并抑制共模干扰，从而为后续放大链路提供稳定、低噪声的正弦基准信号。

\subsection{双级功率放大电路设计}

经 AD8276 转换后的激励信号幅值有限，难以直接驱动 $100\,\mathrm{kHz}$ 条件下具有明显感抗的阵列线圈。因此，本文采用``前置电压放大 + 末级功率驱动''的双级级联拓扑，以兼顾信号保真度、闭环带宽和负载驱动能力。根据硬件设计，AD8276 后级依次采用 OPA197 和 OPA593 两级放大，最终输出激励信号至信号切换板。

\textbf{(1) 前置电压放大级。}

前置电压放大级采用 OPA197 精密运算放大器构建。该级一方面对 AD8276 输出的小幅值正弦信号进行电压放大，另一方面隔离前级 DDS 与末级功率放大器，降低后级负载变化对信号源的影响。通过合理设置闭环增益，前置级可将基准激励信号提升至适合末级功率放大的电平，同时保持较好的幅相稳定性。

\textbf{(2) 末级功率驱动级。}

末级采用 TI 公司的高压大电流功率运算放大器 OPA593。该器件具有较高输出电流能力和较宽供电范围，适合驱动感性线圈负载。在本文系统中，OPA593 主要承担激励链路功率输出任务，将前级已放大的正弦信号进一步提升至约 $\pm12.6\,\mathrm{V}$，并经信号切换板送入被选中的激励线圈。

从系统结构上看，双级级联放大的核心思路是：前级负责高保真电压增益，后级负责负载驱动。与单级大倍率功率放大相比，该结构能够降低末级功放闭环增益压力，改善带宽和稳定性，并降低高频感性负载下的自激风险。

为验证激励链路设计的合理性，以及感性负载驱动能力，本文基于 LTspice 对图 \ref{fig:DDS_AMP} 所示电路进行了仿真分析，其负载按照线圈在100kHz的阻抗实测值进行建模。仿真结果如图 \ref{fig:DDS_AMP_results} 所示，在 $100\,\mathrm{kHz}$ 输入条件下，输出端能够稳定获得约$\pm11.67\,\mathrm{V}$的激励电压；在当前参数设置下，附加相位滞后约为 $14.79^\circ$，总谐波失真低于 $0.126\%$。结果表明，该激励链路能够满足本文 EMT 系统在 $100\,\mathrm{kHz}$ 工作点下的激励需求。

\begin{figure}[!htb]
    \centering
    \includegraphics[width=0.92\textwidth]{figures/DDS_AMP.jpg}
    \figcaption{激励信号发生与双级级联功率放大电路原理图}{Schematic of signal generation and dual-stage cascaded power amplification circuit}
    \label{fig:DDS_AMP}
\end{figure}

\begin{figure}[!htb]
    \centering
    \includegraphics[width=0.92\textwidth]{figures/DDS_AMP_DATA.jpg}
    \figcaption{激励信号发生与双级级联功率放大电路仿真结果}{Simulation results of signal generation and dual-stage cascaded power amplification circuit}
    \label{fig:DDS_AMP_results}
\end{figure}

\section{信号调理与双相锁相放大器设计}

接收线圈拾取到的感应电压信号幅值较小，且容易受到一次背景场、通道串扰、环境电磁噪声以及后级采集链路输入阻抗的影响。若直接将该交流信号送入 ADC 进行数字化采样，不仅会显著提高采样率和动态范围要求，也不利于便携式系统的低功耗设计。因此，本文在接收端设计了一套模拟信号调理与双相锁相放大链路，用于在模拟域完成微弱交流信号的缓冲、增益调节、相敏解调和低通滤波。

该模块的整体功能可概括为三部分：第一，通过硬件自校准、前级高阻缓冲和程控变增益放大，提高接收信号进入锁相解调前的幅值适配能力；第二，通过板级正交 PLL 从激励链路恢复出稳定的同频正交参考信号，为双通道相敏检波提供参考基准；第三，利用两片 AD630 和后级低通滤波网络，将 $100\,\mathrm{kHz}$ 交流响应解调为同相与正交两路直流量，供后级高精度 ADC 采集。

\subsection{硬件自校准机制与前级程控放大电路}
\label{subsec:frontend_vga}

\textbf{1. 硬件自校准机制}

由于 EMT 系统需要在多通道条件下测量微弱信号，各通道开关导通电阻、模拟复用器插入损耗、运算放大器增益误差和电阻分压误差均可能引入系统性偏差。若仅依靠器件标称参数进行补偿，难以满足后续定量成像和模型推理对通道一致性的要求。为此，本文在接收信号链路中加入硬件自校准通道。

具体而言，在接收端模拟复用器与主放大链路之间，系统串联一颗低漏电流、低电荷注入的单刀双掷模拟开关 ADG1419。该开关的常闭端接入来自线圈阵列的实际测量信号，常开端接入一条由激励源分出并经精密衰减得到的同源标准交流信号。该标准信号与系统激励同频同源，经差分恢复和比例衰减后形成约 $5\,\mathrm{mV}$ 的交流基准，用于对后续增益链路进行在线标定。

在系统校准过程中，主控单元将 ADG1419 切换至校准通道，并在不同增益控制电压下记录 ADC 输出结果，从而建立控制电压与系统实际总增益之间的对应关系。该过程将前级开关、缓冲器、程控放大器、相敏检波和低通滤波等链路误差视为整体黑箱进行标定，可有效削弱器件温漂、电阻公差和通道插入损耗对测量结果的影响。

\textbf{2. 前级高输入阻抗缓冲}

从信号切换板输出的接收信号首先进入由 OPA192 构成的高输入阻抗缓冲级。设置该缓冲级的原因在于：阵列线圈和前级选通网络的等效输出阻抗较高，而后级 AD603 程控放大器输入阻抗较低，若二者直接连接，容易形成明显分压，导致微弱感应信号在进入放大器之前已经发生幅值衰减。

OPA192 缓冲级一方面为前端线圈和开关网络提供高输入阻抗负载，降低信号源被后级加载的风险；另一方面以低输出阻抗驱动 AD603 输入端，提高后续程控放大链路的稳定性。该结构能够在不显著引入附加噪声的前提下，改善前后级阻抗匹配关系，为微弱信号增益调节提供稳定输入。

\textbf{3. 程控变增益放大网络}

由于不同线圈对之间的互感响应幅值差异较大，不同伤损位置和尺寸引起的电磁扰动也存在较宽动态范围，固定增益放大结构难以同时兼顾强耦合通道不过载和弱耦合通道可分辨。因此，本文采用 AD603 作为程控变增益放大器，实现接收链路的动态范围调节。

为提高增益控制的稳定性，系统未直接使用 MCU 内部 DAC，而采用独立的 I2C 接口 DAC7571 产生控制电压。DAC7571 输出首先经精密分压网络映射至 AD603 的正向增益控制端，同时利用 REF5025 提供的 $2.5\,\mathrm{V}$ 精密基准经同类型分压网络生成固定偏置，接入 AD603 的负向控制端。由此，AD603 的有效控制电压

\begin{equation}
\label{eq:ad603_control_voltage}
\Delta V_G=V_{\mathrm{GPOS}}-V_{\mathrm{GNEG}}
\end{equation}

被限制在线性控制区间内，避免 DAC 端点非线性和数字电源噪声对增益控制造成不利影响。另外，本文还在系统稳定运行一定时间后测量了AD603增益与VPOS-VNEG电压差之间的映射关系作为进一步校准依据，以提高增益控制的准确性和重复性。

该程控变增益链路可实现约 $10\,\mathrm{dB}$ 至 $50\,\mathrm{dB}$ 的对数增益调节。结合前述硬件校准表和主控端自动增益控制策略，系统能够根据不同通道信号幅值自动选择合适增益，使进入相敏检波器的待解调信号既保持足够幅值，又避免出现饱和失真。程控增益放大及校准链路原理图如图 \ref{fig:AD603} 所示。

\begin{figure}[!htbp]
    \centering
    \includegraphics[width=0.82\textwidth]{figures/AD603.jpg}
    \figcaption{程控变增益放大及校准链路原理图}{Schematic of programmable gain amplification and calibration link}
    \label{fig:AD603}
\end{figure}

\subsection{高速参考信号整形与锁相环控制系统设计}
\label{subsec:pll_design}

双相锁相放大器能否准确提取同相与正交分量，很大程度上取决于参考信号的频率稳定性、相位稳定性和正交精度。若参考信号存在明显相位抖动或正交误差，则解调得到的直流量会直接受到影响。为此，本文设计了一套板级正交 PLL，用于从激励链路中恢复出与测量信号同源的两路 $100\,\mathrm{kHz}$ 正交参考方波。

该板级正交 PLL 由 TLV3603、SN74LV4046 和 SN74ACT74DR 共同构成。其中，TLV3603 用于将激励正弦信号整形成数字方波，SN74LV4046 用于实现鉴频鉴相和压控振荡，SN74ACT74DR 则用于 4 分频和正交参考输出。其控制系统等效结构如图 \ref{fig:PLL_Control_System} 所示。

\begin{figure}[!htbp]
    \centering
    \includegraphics[width=0.78\textwidth]{figures/PLL_Control_System.jpg}
    \figcaption{锁相环控制系统原理图}{Principle diagram of phase-locked loop control system}
    \label{fig:PLL_Control_System}
\end{figure}

\textbf{1. 高速比较器整形与正交参考生成}

来自激励链路的参考正弦信号首先送入 TLV3603 高速比较器。比较器将正弦波过零点转换为数字跳变沿，从而得到与激励信号同频的方波基准。相比直接使用模拟正弦参考，方波参考更适合后级数字锁相环和 AD630 相敏检波器的开关式工作模式。TLV3603 较低的传播延迟和较快的边沿整形能力，可有效降低正弦波过零点附近幅值噪声对时序抖动的影响。

经 TLV3603 整形后的 $100\,\mathrm{kHz}$ 方波被送入 SN74LV4046 的相位比较器。锁定后，VCO 输出约 $400\,\mathrm{kHz}$ 的方波信号，再由 SN74ACT74DR 构成的二阶约翰逊计数器进行 4 分频，得到两路 $100\,\mathrm{kHz}$ 参考方波。其中一路作为反馈信号送回 SN74LV4046，与输入基准进行相位比较；另一路与其保持稳定的 $90^\circ$ 相位差，用于正交通道相敏检波。正交参考生成电路连接拓扑如图 \ref{fig:Second_Order_Johnson_Trigger} 所示。

\begin{figure}[!htbp]
    \centering
    \includegraphics[width=0.78\textwidth]{figures/second_Jonson_trigger.jpg}
    \figcaption{正交参考生成电路连接拓扑}{Connection topology of quadrature reference generation circuit}
    \label{fig:Second_Order_Johnson_Trigger}
\end{figure}

根据 LTspice 仿真结果，TLV3603 引入的附加相位延迟约为 $0.077^\circ$，二阶约翰逊计数器输出的两路 $100\,\mathrm{kHz}$ 方波能够保持稳定正交关系。因此，该结构能够满足双相相敏检波对参考相位一致性的要求。

\textbf{2. 锁相环相位域模型}

为对 PLL 的动态特性进行分析，可将其等效为由鉴相器、环路滤波器、压控振荡器和反馈分频器组成的闭环控制系统。设输入参考相位为 $\Phi_{\mathrm{ref}}(s)$，VCO 输出相位为 $\Phi_{\mathrm{vco}}(s)$，反馈分频比为 $N$，则鉴相器输入相位误差为
\begin{equation}
\label{eq:pll_phase_error}
\Phi_e(s)
=
\Phi_{\mathrm{ref}}(s)
-
\frac{1}{N}\Phi_{\mathrm{vco}}(s)
\end{equation}

SN74LV4046 选用相位比较器 2 作为鉴频鉴相环节。该比较器对边沿敏感，并且对输入方波占空比不敏感，适合本文参考方波链路。在 $5\,\mathrm{V}$ 供电条件下，其小信号鉴相增益可近似为
\begin{equation}
\label{eq:kpd}
K_{PD}
=
\frac{V_{CC}}{4\pi}
\approx
0.398\,\mathrm{V/rad}
\end{equation}

鉴相器输出控制电压与相位误差之间满足
\begin{equation}
\label{eq:vd}
V_d(s)=K_{PD}\Phi_e(s)
\end{equation}

VCO 可视为从控制电压到输出相位的积分环节，其关系为
\begin{equation}
\label{eq:vco_phase}
\Phi_{\mathrm{vco}}(s)
=
\frac{K_{VCO}}{s}V_c(s)
\end{equation}

其中，$K_{VCO}$ 为压控振荡角频率增益。因此，PLL 开环传递函数可写为
\begin{equation}
\label{eq:pll_openloop}
L(s)
=
\frac{K_{PD}K_{VCO}}{N}
\frac{F(s)}{s}
\end{equation}

式中，$F(s)$ 为环路滤波器传递函数。由式 \eqref{eq:pll_openloop} 可知，在 $K_{PD}$、$K_{VCO}$ 和 $N$ 确定后，环路动态性能主要由滤波器 $F(s)$ 决定。

\textbf{3. VCO 参数整定与压控增益测定}

为使 VCO 中心频率落在目标倍频点附近，本文通过硬件调试确定外接电阻、电容参数。在最终配置中，VCO 自由振荡中心频率约为 $422\,\mathrm{kHz}$，能够覆盖系统锁定到 $400\,\mathrm{kHz}$ 所需的牵引范围。

为了获得实际压控增益 $K_{VCO}$，本文对 VCO 控制电压与输出频率之间的关系进行了测量。通过改变 VCO 控制端电压并记录输出频率，可得到图 \ref{fig:K_VCO} 所示的特性曲线。由拟合结果可知，在近似线性工作区间内，VCO 频率斜率约为 $165.27\,\mathrm{kHz/V}$。将其转换为角频率形式，可得
\begin{equation}
\label{eq:kvco}
K_{VCO}
=
2\pi\times165.27\times10^3
\approx
1.038\times10^6\,\mathrm{rad/(s\cdot V)}
\end{equation}

\begin{figure}[!htbp]
    \centering
    \includegraphics[width=0.78\textwidth]{figures/K_VCO.jpg}
    \figcaption{VCO 控制电压与输出频率特性曲线测量及线性拟合}{Characteristic curves of VCO control voltage and output frequency with linear fitting}
    \label{fig:K_VCO}
\end{figure}

由于正交参考生成结构采用 4 分频反馈，因此反馈分频比为
\begin{equation}
\label{eq:pll_n}
N=4.
\end{equation}

结合式 \eqref{eq:kpd} 和式 \eqref{eq:kvco}，可得环路等效增益为
\begin{equation}
\label{eq:pll_loop_gain}
K_L
=
\frac{K_{PD}K_{VCO}}{N}
\approx
1.03\times10^5\,\mathrm{s^{-1}}
\end{equation}

\textbf{4. 环路滤波器设计}

锁相环滤波器的设计目标主要包括三点：第一，闭环带宽应明显低于 $100\,\mathrm{kHz}$ 参考频率，以抑制鉴相器输出中的高频纹波和开关杂散；第二，相位裕度应保持在合理范围内，以避免参考信号扰动下出现明显超调；第三，响应速度应满足阵列通道切换后的稳定时间要求。综合上述要求，本文采用带阻尼补偿和高频旁路的无源 lead-lag 型环路滤波器。

设主积分电阻为 $R_3$，主积分电容为 $C_2$，阻尼补偿电阻为 $R_c$，高频旁路电容为 $C_p$，则环路滤波器可近似表示为
\begin{equation}
\label{eq:pll_filter}
F(s)
=
\frac{1+sR_cC_2}
{
\left[1+s(R_3+R_c)C_2\right]
\left(1+sR_cC_p\right)
}
\end{equation}

其对应零极点频率近似为
\begin{equation}
\label{eq:pll_fp1}
f_{p1}
\approx
\frac{1}{2\pi(R_3+R_c)C_2}
\end{equation}
\begin{equation}
\label{eq:pll_fz}
f_z
\approx
\frac{1}{2\pi R_cC_2}
\end{equation}
\begin{equation}
\label{eq:pll_fp2}
f_{p2}
\approx
\frac{1}{2\pi R_cC_p}
\end{equation}

根据目标穿越频率和相位裕度要求，最终选取
\begin{equation}
\label{eq:pll_component_values}
R_3=10\,\mathrm{k}\Omega,\quad
C_2=1\,\mu\mathrm{F},\quad
R_c=1.2\,\mathrm{k}\Omega,\quad
C_p=47\,\mathrm{nF}
\end{equation}

代入式 \eqref{eq:pll_fp1}--式 \eqref{eq:pll_fp2} 可得
\begin{equation}
\label{eq:pll_freq_values}
f_{p1}\approx14.2\,\mathrm{Hz},\quad
f_z\approx132.6\,\mathrm{Hz},\quad
f_{p2}\approx2.82\,\mathrm{kHz}
\end{equation}

由上述零极点配置可见，阻尼零点位于目标穿越频率以下，用于补偿 VCO 积分环节引入的相位滞后；高频极点则用于进一步衰减鉴相器输出纹波和参考杂散。将式 \eqref{eq:pll_filter} 代入式 \eqref{eq:pll_openloop} 后，可得到完整开环模型，其 Bode 分析结果如图 \ref{fig:PLL_Loop_Filter} 所示。

\begin{figure}[!htb]
    \centering
    \includegraphics[width=0.78\textwidth]{figures/PLL_bode.jpg}
    \figcaption{PLL 环路开环与闭环频率响应}{Open-loop and closed-loop frequency responses of the PLL}
    \label{fig:PLL_Loop_Filter}
\end{figure}

仿真结果表明，该参数组合下系统开环穿越频率约为 $1.54\,\mathrm{kHz}$，闭环 $-3\,\mathrm{dB}$ 带宽约为 $2.48\,\mathrm{kHz}$，相位裕度约为 $60.4^\circ$。该带宽远低于 $100\,\mathrm{kHz}$ 工作频率，能够有效抑制鉴相器输出纹波；同时，约 $60^\circ$ 的相位裕度保证了系统具有较好的阻尼特性和较小超调风险。因此，该 PLL 设计更偏向低噪声、高稳定的正交参考生成，而非追求快速频率跳变跟踪，符合 EMT 微弱信号锁相放大的测量需求。

\subsection{精密相敏检波与低通滤波器设计}
\label{subsec:lockin_design}

锁相放大器的基本原理是将输入信号与已知频率和相位的参考信号进行乘法混频，再通过低通滤波提取直流分量。为便于理解，首先考虑理想情形：设输入信号为 $V_{\mathrm{sig}}(t) = A_s \cos(\omega_0 t + \phi_s)$，同相和正交参考信号分别为正弦波 $\cos(\omega_0 t)$ 和 $\sin(\omega_0 t)$，则相敏检波输出为
\begin{align}
\label{eqn:lockin_I_sine}
V_I &= \frac{2}{T} \int_{0}^{T} V_{\mathrm{sig}}(t) \cos(\omega_0 t) \,\mathrm{d}t
     = A_s \cos\phi_s  \\
\label{eqn:lockin_Q_sine}
V_Q &= \frac{2}{T} \int_{0}^{T} V_{\mathrm{sig}}(t) \sin(\omega_0 t) \,\mathrm{d}t
     = A_s \sin\phi_s 
\end{align}

其中 $T = 2\pi/\omega_0$ 为信号周期。由此可得响应幅值 $A_s = \sqrt{V_I^2 + V_Q^2}$ 和相位\\ $\phi_s = \arctan(V_Q/V_I)$。该过程在频域上等效于以 $\omega_0$ 为中心频率的极窄带滤波，能够从强噪声背景中提取微弱信号。

然而，本文采用的 AD630 为开关型相敏检波器，其参考输入端接收的是由 DDS 产生的 $100\,\mathrm{kHz}$ 正交方波信号，而非上述理想正弦参考。方波与正弦参考的混频机理存在本质差异，需重新推导。设同相参考方波 $R_I(t)$ 和正交参考方波 $R_Q(t)$ 可分别展开为傅里叶级数
\begin{equation}
\label{eq:square_wave_I}
R_I(t) = \frac{4}{\pi} \sum_{n=0}^{\infty} \frac{(-1)^n}{2n+1} \cos\!\bigl((2n+1)\omega_0 t\bigr)
= \frac{4}{\pi}\!\left( \cos\omega_0 t - \frac{1}{3}\cos3\omega_0 t + \cdots \right) 
\end{equation}
\begin{equation}
\label{eq:square_wave_Q}
R_Q(t) = \frac{4}{\pi} \sum_{n=0}^{\infty} \frac{1}{2n+1} \sin\!\bigl((2n+1)\omega_0 t\bigr)
= \frac{4}{\pi}\!\left( \sin\omega_0 t + \frac{1}{3}\sin3\omega_0 t +  \cdots \right) 
\end{equation}

设前级放大后的单频测量信号为
\begin{equation}
\label{eq:sig_signal}
V_{\mathrm{sig}}(t) = V_s \sin(\omega_0 t + \theta) 
\end{equation}

其中 $V_s$ 为信号幅值，$\theta$ 为相对于参考信号的相位差。将该信号分别与同相方波参考 $R_I(t)$ 相乘并经过低通滤波，仅保留直流分量。由于低通滤波器带宽远小于 $\omega_0$，方波中三次及以上谐波分量与基频信号混频产生的高频项均被滤除，仅有基频分量 $\frac{4}{\pi}\cos\omega_0 t$ 参与形成直流输出：
\begin{equation}
\label{eq:square_mix_I}
\begin{split}
V_{\mathrm{sig}}(t) \cdot R_I(t) \big|_{\mathrm{LPF}} 
&= V_s \sin(\omega_0 t + \theta) \cdot \frac{4}{\pi}\cos\omega_0 t \big|_{\mathrm{LPF}} \\
&= \frac{2V_s}{\pi} \sin\theta 
\end{split}
\end{equation}

同理，正交通道经低通滤波后的直流输出为
\begin{equation}
\label{eq:square_mix_Q}
\begin{split}
V_{\mathrm{sig}}(t) \cdot R_Q(t) \big|_{\mathrm{LPF}}
&= V_s \sin(\omega_0 t + \theta) \cdot \frac{4}{\pi}\sin\omega_0 t \big|_{\mathrm{LPF}} \\
&= \frac{2V_s}{\pi} \cos\theta 
\end{split}
\end{equation}

由此可见，在方波参考条件下，同相和正交通道的直流输出分别为
\begin{equation}
\label{eq:lockin_x}
V_X = \frac{2V_s}{\pi} \sin\theta 
\end{equation}
\begin{equation}
\label{eq:lockin_y}
V_Y = \frac{2V_s}{\pi} \cos\theta 
\end{equation}

与正弦参考情形（式 \eqref{eqn:lockin_I_sine}--\eqref{eqn:lockin_Q_sine}）相比，方波参考下的输出幅值比例系数由 $1/2$ 变为 $2/\pi$，其基频分量仍主导解调输出。

通过采集 $V_X$ 和 $V_Y$，系统即可恢复接收通道的复数电压响应：
\begin{equation}
\label{eq:complex_voltage}
V_{\mathrm{meas}} = V_X + \mathrm{j}V_Y 
\end{equation}

其幅值与相位分别为 $\frac{2V_s}{\pi}$ 和 $\arctan(V_X/V_Y)$。该复数响应正是后续传统 EMT 成像和 SCF-DeepONet 模型推理所需的基础观测量。

在硬件连接上，前级放大后的测量信号同时送入两片 AD630 的信号输入端；REF0 和 REF90 两路正交参考方波分别送入两片 AD630 的参考控制端，从而得到同相和正交两路解调输出。AD630 功能结构如图 \ref{fig:AD630_Functional_Block} 所示。

\begin{figure}[!htbp]
    \centering
    \includegraphics[width=0.56\textwidth]{figures/ad630_functional_block.jpg}
    \figcaption{AD630 功能结构图}{Functional block diagram of AD630}
    \label{fig:AD630_Functional_Block}
\end{figure}

\textbf{2. 低通滤波器设计}

相敏检波输出中不仅包含与目标信号相位相关的直流分量，还包含 $100\,\mathrm{kHz}$ 及 $200\,\mathrm{kHz}$ 附近的高频混频分量。为获得稳定直流输出，本文在 AD630 后级设置多阶低通滤波网络，对高频分量进行强抑制，同时保留阵列轮询测量所需的响应速度。

低通滤波器的设计需要在噪声抑制能力和建立时间之间折中。若截止频率过低，虽然能够获得更强的高频抑制能力，但会显著增加通道切换后的等待时间，降低整机单帧采集速度；若截止频率过高，则难以充分滤除混频残余和参考杂散。结合本文阵列轮询测量方式，滤波器被设计为偏低带宽、高抑制的模拟低通网络，以满足微弱信号稳定提取的需求。

\begin{figure}[!htb]
    \centering
    \includegraphics[width=0.78\textwidth]{figures/filter_bode.jpg}
    \figcaption{低通滤波网络频域响应仿真}{AC sweep simulation of the low-pass filter network}
    \label{fig:Filter_Bode}
\end{figure}

\begin{figure}[!htb]
    \centering
    \includegraphics[width=0.78\textwidth]{figures/data_Lock_in_amp.jpg}
    \figcaption{相敏检波器时域响应仿真}{Transient simulation of the phase-sensitive demodulator}
    \label{fig:Lock_in_Amp_Transient}
\end{figure}

为验证滤波网络性能，本文基于 LTspice 进行了频域和时域仿真。频域响应结果如图 \ref{fig:Filter_Bode} 所示。$100\,\mathrm{kHz}$ 处，该滤波器对基频混频分量的抑制达到约 $-154.3\,\mathrm{dB}$；在 $200\,\mathrm{kHz}$ 处，对二倍频分量的抑制约为 $-148.6\,\mathrm{dB}$。该结果表明，所设计的低通网络能够有效削弱相敏检波后的高频残余分量，为后级 ADC 提供较纯净的直流输入。

时域仿真结果如图 \ref{fig:Lock_in_Amp_Transient} 所示。仿真中，相敏检波器启动后，同相与正交通道输出能够平滑过渡至稳态。综合幅值的峰值时间约为 $0.95\,\mathrm{ms}$，超调量约为 $11.9\%$，达到 98\% 稳态所需时间约为 $1.719\,\mathrm{ms}$。稳态综合幅值为 $1.272\,\mathrm{V}$，与理论值 $1.273\,\mathrm{V}$ 的相对误差约为 $0.08\%$。该结果说明，所设计的低通滤波器在保证高频抑制能力的同时，仍能够满足阵列通道轮询测量对建立时间的要求。

综上，本节所设计的信号调理与双相锁相放大链路，实现了从微弱交流感应信号到同相、正交直流量的高信噪比转换。前级硬件校准与程控增益放大提高了通道适配能力，板级正交 PLL 提供了低噪声、相位稳定的参考信号，AD630 与低通滤波网络则完成了窄带相敏解调。该链路为后续高精度模数转换、复数响应构造以及 SCF-DeepONet 模型推理提供了可靠的模拟前端基础。

\section{模数转换与便携式通信模块设计}

经过前级双相正交锁相解调与低通滤波后，接收链路输出已由 100\,kHz 交流响应转换为两路表征电磁扰动信息的低频或直流量，分别对应同相分量与正交分量。为了将这两路模拟量稳定、准确地转换为数字信号，并进一步实现便携式设备与上位机之间的数据交互，本文在测控板后端设计了由电平平移与前端保护、24 位高分辨率模数转换以及主控与蓝牙通信组成的数据采集与传输模块。该模块的总体目标在于：在保证微弱信号量化精度的前提下，实现阵列多通道测量数据的实时组织与无线回传。

\subsection{基于伪差分结构的高分辨率采集电路}

考虑到 EMT 系统中微小伤损往往仅引起极弱的边界电压变化，因此模数转换器不仅需要具有较高分辨率，还需要具备较低本底噪声和较好的低频稳定性。基于这一要求，本设计选用 TI 公司的 24 位 $\Delta$-$\Sigma$ 型模数转换器 ADS1220 作为后端采集核心。将锁相解调后并经过精密电平位移电路输出的两路模拟量分别接入 ADS1220 的两路差分输入，完成同相和正交通道的数字化采样。

如图 \ref{fig:ADS1220_Circuit} 所示，在结构上，本文采用伪差分比例测量思路组织 ADC 输入。具体而言，经过电平平移后的待测信号送入 ADC 的正向输入端，而其负向输入端则绑定至由电压基准芯片 REF5025 提供的稳定参考电位。同时，该高精度 2.5\,V 基准也作为 ADC 的参考端电压。如此，ADC 实际量化的是待测输入相对于基准电平的差值，从而在硬件层面等效扣除了前级平移所引入的中心偏置，并使系统对基准源漂移表现出更高的一致性。

在工作配置上，由于前级已经设置了程控变增益放大器完成动态范围匹配，因此 ADS1220 内部 PGA 可不作为主要增益来源使用，从而避免引入额外低频噪声与失调漂移。与此同时，系统可根据测量速度与噪声水平之间的折中关系，对 ADS1220 的输出速率进行配置，以在阵列分时扫描中兼顾单次采样精度和整轮扫描时间。

综合来看，该采集架构的特点在于：一方面，利用锁相解调将高频交流信息前移至模拟域完成压缩，使 ADC 只需面对低带宽、高精度直流量；另一方面，通过伪差分测量和统一参考基准设计，提高了双通道 I/Q 采样的一致性，为后续复数响应构造提供了可靠的数据基础。

\begin{figure}[!htbp]
    \centering
    \includegraphics[width=0.82\textwidth]{figures/ADC.jpg}
    \figcaption{高精度 ADC 采集电路}{High-precision ADC acquisition circuit}
    \label{fig:ADS1220_Circuit}
\end{figure}

\subsection{主控单元与蓝牙通信协议设计}

本系统选用基于 ARM Cortex-M4 内核的 STM32F407VGT6 作为主控单元。该 MCU 主要承担以下任务：激励源配置、继电器与模拟开关时序控制、DAC 增益设置、校准流程调度、ADC 数据读取以及测量结果封包与上传。对于本文所设计的便携式 EMT 平台而言，主控的核心作用并不在于完成复杂成像计算，而在于以稳定时序驱动整机完成多通道循环测量，并将整理后的数据可靠传输至上位机。

为了满足现场检测中对便携性和去有线连接化的需求，系统在主控板上集成了蓝牙无线透传链路，用于实现测量结果的实时发送。下位机与上位机之间采用轻量级自定义数据帧协议进行通信。上位机可向下位机发送启动测量、停止测量、增益校准、空场基准校准等控制命令；而下位机则在每轮完整扫描结束后，将当前测量状态、通道信息以及采集数据组织成固定格式的数据帧后上传。

需要特别说明的是，对于 4$\times$4 阵列 EMT 系统而言，一轮完整扫描并非只产生少量通道数据，而是对应 16 次激励工况下、每次 15 个接收通道的完整测量结果。因而，下位机上传的核心数据应组织为一组完整的多视角复数响应观测值，即包含同相与正交信息的 240 组测量结果。必要时，还可附加激励电流监测结果、设备状态字和校验字段，以保证后续数据解析与异常诊断的可靠性。这样组织后的数据既可直接用于第二章中的传统 EMT 成像算法，也可进一步送入第三、四章所提出的 SCF-DeepONet 模型进行伤损热力图推理。

这种“前端模拟窄带解调 + 后端低速高精度采样 + 无线数字回传”的架构具有明显优势：它将高频模拟信号处理尽可能前移至板级模拟域完成，大幅降低了数字侧对高速采样和大吞吐传输的依赖；同时，将复杂的图像重建计算任务保留给上位机执行，使得便携式下位机能够在较低功耗和较小体积条件下完成稳定的数据采集与传输任务。

\section{系统电源管理与启机时序设计}

本系统同时集成了大电流功率放大、高增益微弱信号调理、高精度模数转换以及数字主控与无线通信等多类功能模块。不同模块在供电电压等级、瞬态电流需求以及电源噪声容限方面存在显著差异：功率驱动级要求较高电压和较强瞬态供电能力，锁相解调与 ADC 前端则对纹波、地噪声及低频漂移极为敏感，而 MCU 与数字逻辑电路又会在工作过程中引入明显的开关噪声。为降低各模块之间的电源耦合干扰，本文构建了“DC-DC 预稳压 + LDO 精密稳压”的分布式低噪声电源树结构，并进一步设计了面向多电源轨模数混合系统的安全启机时序。

如图 \ref{fig:Power_Tree} 所示，系统以外部直流输入作为主电源入口，并根据不同功能模块的供电需求，采用“一级高效变换、二级低噪声稳压”的方式划分多条独立电源轨。其设计目标是：先通过开关电源完成电压粗调和极性转换，再通过线性稳压器对关键模拟链路进行低噪声供电，从而兼顾转换效率与模拟精度。

\begin{figure}[!htbp]
    \centering
    \includegraphics[width=0.82\textwidth]{figures/power_tree.jpg}
    \figcaption{分布式低噪声电源树拓扑}{Distributed low-noise power tree topology}
    \label{fig:Power_Tree}
\end{figure}

对于 DDS、ADC、DAC 以及 MCU 等模数混合器件，系统进一步采用局部独立供电策略。根据当前硬件方案，AD9834 与 ADS1220 的模拟供电由 5\,V 支路提供，AD603 则采用 $\pm 5$\,V 供电，主控 STM32F407 及其数字逻辑接口工作于 3.3\,V 电源域。各支路均尽量采用就近稳压、就近去耦的方式布置，以减少数字侧高频开关电流对模拟基准和敏感采样链路的污染。同时，模拟供电与数字供电在 PCB 上采取分区域布设和局部隔离策略，从供电路径上降低不同噪声源之间的直接耦合。

综合来看，该分布式电源树并非单纯追求“电压够用”，而是围绕模数混合系统中的噪声传播路径进行针对性设计：高功率链路优先保证供电能力，微弱信号链路优先保证低纹波和低漂移，数字链路则优先考虑稳定驱动与时序可靠性。通过电源树分级和功能分域设计，为后续高精度测量提供了基础保障。

\section{系统硬件整体布线与 PCB 抗干扰设计}

本系统同时包含 100\,kHz 大电流信号激励链路、微弱信号锁相调理链路、高速数字逻辑以及开关电源模块，属于典型的高集成模数混合系统。对于此类系统而言，硬件性能不仅取决于原理图设计，还高度依赖 PCB 层面的布局布线、电源回流路径控制以及地划分策略。若这些细节处理不当，即使各单元电路在仿真或单板测试中性能良好，整机联调后仍可能出现底噪抬升、相位漂移、串扰增大甚至测量重复性下降等问题。因此，本文在空间结构、板层规划、敏感信号布线以及接地与屏蔽策略等方面均进行了针对性抗干扰设计。

\subsection{空间堆叠架构与四层板功能分区布局}

在系统总体硬件结构上，本文采用``信号切换板 + 测控主板''分离的双板式方案。两块板之间通过FPC连接器进行数字信号的通信，通过屏蔽良好的SMA连接完成关键信号交互，并通过铜螺柱进行机械固定。这种空间分离结构有利于将大电流切换和继电器动作所引起的瞬态干扰，与高增益模拟测量链路在物理上分开布置，从而减小前端强干扰对高精度采集链路的直接耦合。

如图 \ref{fig:mux_board} 为信号切换板的实物照片。该板主要承担继电器选通、模拟开关控制以及部分激励链路的功能。其设计重点在于：将大电流切换相关的继电器驱动与激励链路布设在相对独立的区域，并通过合理的走线和地划分策略降低其对后级测量链路的干扰；同时，确保关键信号路径尽量短、过孔少，并通过伴随地线和局部屏蔽等方式提高其抗干扰能力。

\begin{figure}[htbp]
\centering
% --- 子图 a：正面 ---
\begin{minipage}[t]{0.49\textwidth}
\centering
\includegraphics[width=0.99\textwidth]{figures/21.jpg}\\
{\zihao{5}(a) 正面 \\ {\zihao{5}(a) Front View}}
\end{minipage}
% --- 子图 b：反面 ---
\begin{minipage}[t]{0.49\textwidth}
\centering
\includegraphics[width=0.99\textwidth]{figures/22.jpg}\\
{\zihao{5}(b) 反面 \\ {\zihao{5}(b) Back View}}
\end{minipage}
% --- 总图题 ---
\figcaption{多路复用切换板实物图}{Physical photo of the multiplexer board}
\label{fig:mux_board}
\end{figure}

测控主板的实物图如图 \ref{fig:measure_board} 所示，本设计采用四层 PCB 结构进行实现。顶层用于关键模拟信号和高频敏感走线，第二层布设完整连续的地平面，第三层用于电源分配，底层则主要承载数字逻辑和控制信号。完整的实体地平面不仅能够为高频和瞬态电流提供低阻抗回流路径，也有助于减小电流环路面积，降低空间辐射与互感耦合。

在器件布局上，系统遵循``按信号流线性展开''和``高噪声源边缘化''的原则。开关电源模块、MCU 与数字逻辑、功率放大模块以及高精度锁相放大链路分别布设在不同区域，并尽量按照信号流向依次排列。其中，开关电源和数字时钟链路布置在远离高阻抗模拟输入的位置；锁相放大与 ADC 前端则集中布设在相对独立、干扰更小的区域。这样的功能分区有利于降低不同噪声源之间的近场耦合，提高整板的电磁兼容性。

\begin{figure}[ht]
    \centering
    \includegraphics[width=0.72\textwidth]{figures/11.jpg}
    \figcaption{测控主板实物图}{Physical photo of the measurement board}
    \label{fig:measure_board}
\end{figure}

\subsection{关键信号布线与多地隔离防护策略}

在接收链路中，从信号切换板引入的微弱信号具有幅值低、源阻抗较高、对寄生参数敏感等特点。因此，对于该类关键信号，本文在 PCB 布线中采用了短路径、少过孔、伴随地包围以及局部屏蔽等策略，以尽量降低漏电流、寄生电容和空间辐射引起的附加误差。尤其是在进入 OPA192 缓冲器之前的高阻抗节点附近，应避免长距离平行走线，并通过伴随地线与密集接地过孔提高其抗干扰能力。

对于锁相参考链路、AD630 解调输出以及 ADC 输入等关键模拟节点，系统尽量采用同层短距离连接，减少不必要的跨层切换，以降低附加寄生电感和寄生电容对相位与幅值的一致性影响。与此同时，模拟参考电压、增益控制电压以及 ADC 基准引线均尽量避开高频数字总线和开关节点，避免将数字时序噪声直接注入精密模拟链路。

在接地拓扑方面，本文采用``按功能分区、局部隔离、单点汇接''的设计思想。对于信号切换板，大电流继电器驱动和激励链路相关地回路与微弱接收链路相关模拟地进行了分开处理；对于测控主板，数字地、功率地和敏感模拟地在布局上分别占据不同区域，并在合适位置进行单点连接，以降低大电流动态回流经过模拟测量区域所引起的地弹噪声。需要指出的是，这里隔离并非指绝对断开，而是强调对不同类型回流路径进行约束和管理，使高频脉冲电流与微弱模拟回流电流尽可能不共享关键路径。

此外，对于外部接口、调试连接口以及可能与金属结构件相连的安装点，还考虑了静电与地环路问题。为此，采用了由高阻电阻与小电容构成的 RC 接地网络，在低频条件下抑制地环路电流，而在高频静电或尖峰脉冲出现时为其提供泄放路径。这样既有助于减小外部环境对整机参考地的扰动，也有利于提高系统在现场使用条件下的抗静电能力。

总体而言，本文所采用的双板空间分离、四层板分区布局、关键模拟线伴随包地以及多地单点汇接等措施，共同构成了 EMT 硬件平台的 PCB 抗干扰设计基础。它们并不直接提高系统的理论分辨率，却对实际整机的噪声底、重复性和长期稳定性起到了决定性作用，是便携式高精度 EMT 测量系统实现中不可忽视的重要环节。

\section{本章小结}

本章围绕便携式道岔伤损 EMT 检测需求，完成了系统硬件平台的总体设计与实现。首先，结合高电导率、高磁导率铁磁目标的检测特点，明确了系统在激励能力、微弱信号提取、阵列复用、便携供电、频率可调以及现场抗干扰等方面的核心需求；在此基础上，构建了由``信号切换板 + 测控板''组成的双板式总体硬件架构，并给出了其完整信号流与工作闭环。

随后，本文分别对各关键硬件模块进行了设计分析：针对前端传感器部分，设计并制作了 $4\times4$ 平面线圈阵列，并通过阻抗测试与仿真对比验证了线圈参数一致性；针对阵列分时复用需求，提出了基于机械继电器与 ADG1206 的收发异构多路选通结构，并引入激励电流监测链路；针对激励端，采用 AD9834、AD8276、OPA197 与 OPA593 组成 100\,kHz 主工作点下的激励产生与功率放大链路，同时保留了一定频率调整能力；针对接收端，则设计了由 OPA192 缓冲、AD603 程控增益放大、板级正交 PLL、双 AD630 锁相解调与低通滤波构成的高信噪比模拟调理链路，并通过 ADS1220 与 STM32 完成 I/Q 两路信号的高精度采集、封包与无线回传。

最后，围绕整机实现，本文进一步给出了分布式低噪声电源树、软硬协同启机时序以及双板空间分离、四层板分区布局、关键模拟线伴随包地和多地单点汇接等 PCB 抗干扰设计策略。上述硬件设计共同构成了面向 100\,kHz EMT 伤损检测任务的完整物理平台，为后续系统联调、传统成像算法验证以及 SCF-DeepONet 模型推理提供了可靠的实验基础。
